\documentclass[14pt]{extarticle}
\usepackage{preamble}
\geometry{letterpaper, landscape, margin=0.5in}
\usepackage{qrcode}
\renewcommand{\answer}[1]{}
\setlength{\parskip}{4pt}

\begin{document}
\pagestyle{empty}

\daybanner{Unit 3 \textperiodcentered\ Topic 3.9 \textperiodcentered\ TODAY'S PROBLEM}{Two Ways to Sample Library Visits}

\noindent\begin{minipage}[t]{0.57\linewidth}
\vspace{0pt}
A library studies self-checkout before and after a redesign. Records give $p_B=0.42$ among 4{,}000 before visits and $p_A=0.54$ among 5{,}000 after visits.\par\smallskip \textbf{Plan I:} Draw an SRS of 120 before visits and a separate SRS of 150 after visits; do not track cardholders.\par \textbf{Plan P:} Measure the same 120 cardholders before and after.\par\smallskip \textbf{Sora:} Separate selection warrants Plan I's independent model; 0.0350 means a difference at most 0.01 is unusual.\par \textbf{Mateo:} Pairing controls person differences, so Plan P may help. Because both proportions use 120 people, Plan I's model applies to P too.\par\smallskip For Plan I, the four expected counts are 81, 69, 50.4, and 69.6. A normal calculator gives the lower-tail probability 0.0350.\par
\end{minipage}\hfill
\begin{minipage}[t]{0.40\linewidth}
\vspace{0pt}
\begin{center}
\textbf{Sampling plans}\par\smallskip
\begin{tabular}{|l|c|c|c|}
\hline
\textbf{Plan} & \textbf{Before} & \textbf{After} & \textbf{Repeated} \\ \hline
\textbf{I} & SRS 120 & SRS 150 & no \\ \hline
\textbf{P} & 120 & same 120 & yes \\ \hline
\end{tabular}
\end{center}
\end{minipage}

\begin{firsttakebox}{FIRST TAKE (5 min, on your sheet)}
Does using the same people twice seem different from choosing two separate groups? Give one reason.
\end{firsttakebox}

\begin{description}[leftmargin=1.15in, labelwidth=1in, itemsep=2pt, topsep=2pt]
  \item[\dokbadge{1}~(a)] For after minus before, identify the parameter and statistic and state Plan I's sampling-distribution mean.
  \item[\dokbadge{2}~(b)] For Plan I, verify both 10\% conditions and use the four supplied counts to check normality. Calculate the SD and interpret the supplied probability $P(\widehat p_A-\widehat p_B\leq0.01)=0.0350$.
  \item[\dokbadge{3}~(c)] In 3--4 sentences, explain what Sora and Mateo each get right and wrong. Use how the samples are selected to decide where the independent model applies. Explain why transferring its SD and probability to the repeated people could mislead a reader.
\end{description}

\vfill
\noindent\begin{minipage}{0.84\linewidth}
\textbf{Video + follow-along:} \href{https://robjohncolson.github.io/apstats-live-worksheet/u5_lesson6_live.html}{\texttt{\detokenize{u5_lesson6_live.html}}} (scan the code)\par
\textbf{Finish (a)--(c) after the video. Turn in your sheet.}
\end{minipage}\hfill
\qrcode[height=0.75in]{https://robjohncolson.github.io/apstats-live-worksheet/u5_lesson6_live.html}

\end{document}
